\documentclass[10pt]{article}
%% korr-preamble.tex -- gemeinsame Präambel für die Korrekturlisten LNM 1184
%% Makros wie im Buch (LN-Book-Springer/preamble), damit die Formeln
%% genauso aussehen wie im Proof.
\usepackage[T1]{fontenc}
\usepackage[utf8]{inputenc}
\usepackage{newtxtext}
\usepackage[varvw]{newtxmath}
\usepackage[scaled=.9]{helvet}
\usepackage[varqu,scaled=.95]{zi4}
\usepackage[ngerman]{babel}
\usepackage{csquotes}
\usepackage[a4paper,landscape,margin=1.6cm,top=1.8cm,bottom=1.8cm]{geometry}
\usepackage{microtype}
\usepackage{array,longtable,booktabs}
\usepackage{ragged2e}
\usepackage[table]{xcolor}
\usepackage{xspace}
\usepackage{enumitem}
\setlist{nosep,leftmargin=*}
\usepackage[nointegrals]{wasysym}
\usepackage{newunicodechar}
\usepackage{nicefrac}
\usepackage{mathtools}
\usepackage{mleftright}
\mleftright
\usepackage[hidelinks]{hyperref}
\setlength{\parindent}{0pt}
\setlength{\parskip}{.5em}
\setlength{\emergencystretch}{3em}
\renewcommand{\arraystretch}{1.25}
\definecolor{codegray}{gray}{0.35}

%% ---- Makros aus dem Buch -------------------------------------------
\DeclarePairedDelimiterX{\norm}[1]{\lVert}{\rVert}{\ifblank{#1}{\:\cdot\:}{#1}}
\DeclarePairedDelimiterX{\abs}[1]{\lvert}{\rvert}{\ifblank{#1}{\:\cdot\:}{#1}}
\DeclarePairedDelimiterX{\dualp}[1]{\langle}{\rangle}{\ifblank{#1}{\,\cdot\,,\cdot\,}{#1}}
\newcommand{\interval}[1]{\left[ #1 \right]}
\newcommand{\ointerval}[1]{\left] #1 \right[}
\newcommand{\rointerval}[1]{\left[ #1 \right[}
\newcommand{\lointerval}[1]{\left] #1 \right]}
\newcommand{\tint}{\textstyle\int}
\AtBeginDocument{%
\let\ORGvarepsilon=\varepsilon
\let\epsilon=\epsilon
\let\epsilon=\ORGvarepsilon
\let\ORGvarrho=\varrho
\let\rho=\rho
\let\rho=\ORGvarrho
\let\ORGvartheta=\vartheta
\let\theta=\theta
\let\theta=\ORGvartheta
\let\ORGvarphi=\varphi
\let\phi=\phi
\let\phi=\ORGvarphi
\let\ORGvarleq=\leqslant
\let\leqslant=\leq
\let\leq=\ORGvarleq
\let\ORGvargeq=\geqslant
\let\geqslant=\geq
\let\geq=\ORGvargeq
\let\ORGvarpreccurlyeq=\preccurlyeq
\let\preccurlyeq=\preceq
\let\preceq=\ORGvarpreccurlyeq
\let\setminus\smallsetminus			% A - B fuer Mengen
\renewcommand{\P}{\mathfrak{P}}		% Achtung: \P hat andere Bedeutung in LaTeX
}
\newcommand{\N}{\mathbb{N}}		% Natuerliche Zahlen
\newcommand{\Z}{\mathbb{Z}}		% Ganze Zahlen
\newcommand{\Q}{\mathbb{Q}}		% Rationale Zahlen
\newcommand{\R}{\mathbb{R}}		% Reelle Zahlen
\newcommand{\C}{\mathbb{C}}		% Komplexe Zahlen
\newcommand{\K}{\mathbb{K}}		% Koerperzeichen
\newcommand{\T}{\mathbb{T}}		% Torus
\DeclareMathOperator{\Id}{Id}
\DeclareMathOperator{\sign}{sign}				% signum-Operator
\newcommand{\im}{\ensuremath{\mathrm{i}}}		% imaginäre Konstante i
\newcommand{\iu}{\ensuremath{\mathrm{i}}}
\newcommand{\eu}{\ensuremath{\mathrm{e}}}
\renewcommand*{\Re}{\mathop{}\!\mathrm{Re}}		% Realteil
\renewcommand*{\Im}{\mathop{}\!\mathrm{Im}}		% Imaginärteil
\newcommand*{\supp}{\mathop{}\!\mathrm{supp}\,}  % Träger
\newcommand*{\diff}[1]{\mathop{}\!\mathrm{d}{#1}}	% \diff{s} = ds, 
\newcommand*{\ds}{\mathop{}\!\mathrm{d}{s}}       	% \ds = ds, 
\newcommand*{\dg}{\mathop{}\!\mathrm{d}{g}}       	% \dg = dg, 
\newcommand*{\dt}{\mathop{}\!\mathrm{d}{t}}       	% \dt = dt, 
\newcommand*{\dx}{\mathop{}\!\mathrm{d}{x}}       	% \dx = dx, 
\newcommand*{\dy}{\mathop{}\!\mathrm{d}{y}}       	% \dy = dy, 
\newcommand*{\dr}{\mathop{}\!\mathrm{d}{r}}       	% \dr = dr, 
\newcommand*{\dm}{\mathop{}\!\mathrm{d}{m}}       	% \dm = dm, 
\newcommand*{\du}{\mathop{}\!\mathrm{d}{u}}       	% \du = du, 
\DeclareMathOperator{\esssup}{ess\,sup}
\newcommand{\LE}{\mathcal{L}(E)} 		% 
\newcommand{\BH}{\mathcal{B}(H)} 		% 
\newcommand{\BA}{\mathfrak{A}}
\newcommand{\ZE}{\mathcal{Z}(E)} 	% \ZE  Zentrum von E
\newcommand{\LH}{\mathcal{L}(H)} 
\let\originalL\L
\AtBeginDocument{%
  \DeclareRobustCommand{\L}[1]{\mathcal{L}(#1)}%
}
\newcommand{\Fix}[1]{\mathop{}\mathrm{Fix}\left(#1\right)}		% \Fix{T} Fixraum
\newcommand{\rank}{\mathop{}\!\mathrm{rank}\,}		% 
\newcommand{\Kern}[1]{\mathop{}\mathrm{ker}(#1)}
\newcommand{\Image}[1]{\mathop{}\mathrm{im}(#1)}
\newcommand{\image}{\mathop{}\!\mathrm{im}\,}
\renewcommand{\dim}{\mathop{}\!\mathrm{dim}\,}
\newcommand*{\eg}{e.g.,\xspace}
\newcommand*{\ie}{i.e.,\xspace}
\newcommand*{\resp}{resp.\xspace}
\newcommand*{\vs}{vs.\xspace}
\newcommand*{\cf}{cf.\xspace}
\newcommand*{\etc}{etc.\xspace}
\newcommand{\CA}{$\mathrm{C}^{*}$}	% C*-Algebra
\newcommand{\WA}{$\mathrm{W}^{*}$}	% W*-Algbera
\newcommand{\AW}{$\mathrm{AW}^{*}$}	% AW*-Algbera
\newcommand*{\TT}{\mathcal{T}}		% C_0-Semigroup
\newcommand*{\RR}{\mathcal{R}}		% Pseudo Resolvente
\newcommand*{\SG}{\mathcal{S}}	    % Semigroup
\newcommand*{\UG}{\mathcal{U}}	    % Implemented Semigroup
\newcommand*{\F}{\mathcal{F}}		% F-product
\newcommand{\mfrac}[2]{%
    {\raise.5ex\hbox{$\scriptstyle #1$}}%
    \kern.025em/\kern.025em%
    {\lower.25ex\hbox{$\scriptstyle #2$}}}
\newcommand*{\sfrac}[2]{\leavevmode\kern.1em%
       \raise.5ex\hbox{\small #1}%
       \kern.025em/\kern.025em%
       \lower.25ex\hbox{\small #2}}
\newcommand{\1}{\mathbf{1}}
\newcommand{\HO}{\mathcal{H}(O)} 		%  für Lotz
\newcommand{\WO}{\mathcal{W}(O)} 		%  für Lotz
\newcommand*{\ddp}{\mathop{}\!\mathrm{d}{p}}       	% \ddp = dp, 
\newcommand*{\dN}{\mathop{}\!\mathrm{d}{N}}       	% \dN = dN, 
\newcommand*{\ddt}{\mathop{}\!\mathrm\fraq{d}{dt}}   % \ddt = d/dt
\newcommand*{\lnm}{LNM1184: } % \LMN-BookCode
\newcommand{\closure}[1]{#1^{-}}
%% Hilfsmatrizen aus Teil I, Kap. 1
\newcommand{\ABC}{\begin{pmatrix} 0 & 1 \\ -1 & 0 \end{pmatrix}}
\newcommand{\BBC}{\begin{pmatrix} \cos(t) & \sin(t) \\ - \sin(t) & \cos(t) \end{pmatrix}}
\newcommand{\CBC}{\begin{pmatrix} 1 & 1 \\ -1 & -1 \end{pmatrix}}
\newcommand{\DBC}{\begin{pmatrix} 1+t & t \\ - t & 1-t \end{pmatrix}}
\newcommand{\EBC}{\begin{pmatrix} 0 & 1 \\ 1 &0 \end{pmatrix}}
\newcommand{\FBC}{\begin{pmatrix} \cosh(t) & \sinh(t) \\ \sinh(t) & \cosh(t) \end{pmatrix}}

%% ---- Darstellung der Listen ----------------------------------------
%% Quelltext, der sich nicht als Formel setzen lässt (Labels, \ref, Dateinamen)
\newcommand{\code}[1]{\leavevmode{\color{codegray}\ttfamily\small #1}}
\newcommand{\sichA}{\CIRCLE}      % ● hoch
\newcommand{\sichB}{\LEFTcircle}  % ◐ mittel
\newcommand{\sichC}{\Circle}      % ○ niedrig
%% Querverweise im Buch: Label-Namen statt Nummern
\newcommand{\lbl}[1]{\textsf{\footnotesize\color{codegray}\detokenize{#1}}}

\hypersetup{pdftitle={Formelprüfung Part A (Teil I): Satz und Notation}}
\pagestyle{plain}
\begin{document}
% Erzeugt aus korrekturen-mathe-part-a.md: Formeln gesetzt statt Quelltext.
% Grau/Schreibmaschine = Quelltext, der sich nicht als Formel setzen laesst (Makronamen, Labels).
{\Large\bfseries Formelprüfung Part A (Teil I): Satz und Notation\par}
\smallskip
\emph{One-parameter Semigroups of Positive Operators}, 2. Aufl. (Springer LNM 1184). Stand: 14.10.2026.

\textbf{Umfang:} Teil I, Kap. 1–4, geprüft im LaTeX-Quelltext (Branch \code{fahnenkorrektur}). Geprüft wurden nur Satz und Notation in Formeln: Tippfehler, Klammern, Indizes, uneinheitliche Schreibweisen, Mathemodus, Satzzeichen an Formeln, Gleichungsverweise. Beweise wurden \textbf{nicht} nachgerechnet; wo ein Befund inhaltlich wirkt, steht „bitte prüfen“.

\textbf{Noch nichts geändert.} Die Liste dient zur Besprechung; umgesetzt wird erst nach Rücksprache.

\textbf{Spalten:} S. $=$ gedruckte Buchseite, ermittelt über das nächstgelegene Label (kann um eine Seite abweichen). Z. $=$ Zeile in der .tex-Datei. Sicherheit: \sichA{} eindeutig $\ensuremath{\cdot}$ \sichB{} sehr wahrscheinlich $\ensuremath{\cdot}$ \sichC{} Hinweis/Ermessen.

Insgesamt \textbf{141 Punkte}: 34 \sichA{}, 58 \sichB{}, 49 \sichC{}. Davon stehen 9 bereits in der Arbeitsliste \code{Fahnenkorrektur\_Springer\_LNM1184.md} (als „Schon in der Arbeitsliste“ markiert).

\section*{Auffälligste Punkte}
\begin{itemize}
\item \textbf{S. 73, Example 3.1.4:} Das optionale Argument steht hinter \code{\textbackslash{}label}; „[Hilbert space, Zabczyk]“ wird deshalb als Fließtext gedruckt statt im Kopf des Beispiels.
\item \textbf{Mehrere Klammerfehler in Formeln} (Kap. 2: Z. 657, 1645, 1675, 1679; Kap. 1: Z. 767).
\item \textbf{Vertauschte oder falsche Variablen/Indizes}, z. B. $1/2$ statt $s/2$ (Kap. 1, Z. 470), $f_f'$ statt $\xi_f'$ (Kap. 1, Z. 1332/1337), $\lim_{t\to0}$ statt $\lim_{n\to\infty}$ (Kap. 2, Z. 2031), $t\to\infty$ statt $t\to0$ (Kap. 3, Z. 1155), fehlendes Minus in der Laplace-Transformation (Kap. 4, Z. 205).
\item \textbf{Tilden \code{\textasciitilde{}} im Mathemodus} erzeugen sichtbaren Leerraum (Kap. 2, Z. 1473, 1914, 1951).
\end{itemize}
\section*{Einzelbefunde}
\subsection*{Teil I, Kap. 1 – Basic Results on Semigroups on Banach Spaces (\code{part-a/chap-a1.tex})}
\begin{longtable}{@{}rrc>{\RaggedRight\arraybackslash}p{5.8cm}>{\RaggedRight\arraybackslash}p{5.8cm}>{\RaggedRight\arraybackslash}p{1.9cm}>{\RaggedRight\arraybackslash}p{7.3cm}@{}}
\toprule
\textbf{S.} & \textbf{Z.} & \textbf{} & \textbf{Ist} & \textbf{Soll} & \textbf{Art} & \textbf{Kommentar} \\
\midrule
\endfirsthead
\toprule
\textbf{S.} & \textbf{Z.} & \textbf{} & \textbf{Ist} & \textbf{Soll} & \textbf{Art} & \textbf{Kommentar} \\
\midrule
\endhead
\bottomrule
\endlastfoot
4 & 69 & \sichB{} & $T(t)_{t\geq0}$ & $(T(t))_{t\geq0}$ & Klammern & Sonst überall $(T(t))_{t\geq0}$; ebenso Z. 78. \\[2pt]
4 & 104 & \sichB{} & $-\infty \leq \omega < +\infty$ & $-\infty \leq \omega_{0} < +\infty$ & Tippfehler & Im Satz geht es um $\ensuremath{\omega}\ensuremath{_{0}}$. \\[2pt]
6 & 220 & \sichB{} & $\frac{d}{\dt}\xi(t)$ & $\frac{\diff{}}{\dt}\xi(t)$ & Notation & Kursives d; in Z. 187, 321, 648 aufrecht. Ebenso Z. 785 $\frac{d}{\dx}$. \\[2pt]
7 & 285 & \sichB{} & $\text{ for $ n \geq 2 $, }$ & $\text{ for $ n \geq 2 $.}$ & Satzzeichen & Danach beginnt ein neuer Satz („Then“, Z. 288). \\[2pt]
9 & 380 & \sichB{} & $= \colon  \eu^{tA}$ & $\eqqcolon \eu^{tA}$ & LaTeX & $= \colon$ ergibt „$=$ :“ mit falschem Abstand; ebenso Z. 561 und 818. \\[2pt]
9 & 448 & \sichB{} & {for all $ (k \geq 2)$} & {for all $ k \geq 2 $} & Klammern & Überflüssige Klammern. \\[2pt]
9 & 470 & \sichA{} & $\eu^{t( A - 1/2  \cdot \Id )}$ & $\eu^{t( A - s/2  \cdot \Id )}$ & Tippfehler & In Z. 451 und 471 steht s/2. \\[2pt]
9 & 472 & \sichB{} & $\sum_{ \text{k even} }$ & $\sum_{k \text{ even}}$ & Mathemodus & k aufrecht in Textschrift; ebenso $\text{k odd}$. \\[2pt]
14 & 767 & \sichA{} & \code{\dots{}\$Af \allowbreak{}\allowbreak{}\textbackslash{}coloneq \allowbreak{}f'\$\}. \allowbreak{}\allowbreak{}\textbackslash{}\}} & $\dots , f' \in E\}, \quad Af \coloneq f'.$ & Klammern & Punkt vor $\}$, und „Af := $f\ensuremath{{}'}$“ steht innerhalb der Mengenklammer. \\[2pt]
15 & 814 & \sichC{} & $\cos \pi nx$ & $\cos(\pi nx)$ & Notation & Z. 851 mit Klammern, Z. 814 und 858 ohne. \\[2pt]
15 & 871 & \sichA{} & $k_{t}(x,y)dy$ & $k_{t}(x,y)\dy$ & Notation & Kursives d ohne Abstand; sonst $\dy$. Ebenso Z. 884 $f(y)dy$. \\[2pt]
16 & 878 & \sichB{} & \code{\textbackslash{}subsection\{n-dimensional \allowbreak{}\dots{}\}} & \code{\textbackslash{}subsection\{\allowbreak{}\textbackslash{}texorpdfstring\{\$n\$\}\{n\}-dimensional \allowbreak{}\dots{}\}} & Mathemodus & n nicht im Mathemodus (im Index steht $n$). \\[2pt]
16 & 881 & \sichC{} & $1 \leq p \leq \infty$ & $1 \leq p < \infty$ & Tippfehler & Auf $L^\ensuremath{\infty}$ nicht stark stetig, $\ensuremath{\mathcal{S}}$ dort nicht dicht (Z. 906); inhaltlich bitte prüfen. \\[2pt]
16 & 916 & \sichA{} & $T(t)\vert_{ S(\R^{n})}$ & $T(t)\vert_{\mathcal{S}(\R^{n})}$ & Notation & Schwartz-Raum sonst $\mathcal{S}$. \\[2pt]
16 & 952 & \sichC{} & $L^{p}\interval{0,1}$ & $L^{p}(\interval{0,1})$ & Notation & Z. 755 und 1176: $L^p([0,1])$. \\[2pt]
19 & 1047 & \sichB{} & $(T(t)_{\vert})$ & $(T_{\vert}(t))_{t \geq 0}$ & Notation & Z. 1035 definiert $T_{\vert}(t)$; Einschränkungen insgesamt uneinheitlich notiert (Z. 1167, 1251, 1379). \\[2pt]
20 & 1109 & \sichC{} & $\hat{g} \, \ds$ & $\hat{g} \ds$ & Abstand & Zusätzliches \code{\textbackslash{},} nur in Z. 1109–1111. \\[2pt]
20 & 1115 & \sichC{} & $\lambda > \omega$ & $\lambda > \omega_{0}$ & Notation & $\ensuremath{\omega}$ ohne Index hier nicht definiert. \\[2pt]
20 & 1159 & \sichB{} & {(see Example~\lbl{subsec:a1-2.4})} & {(see Section~\lbl{subsec:a1-2.4})} & Verweis & Schon in der Arbeitsliste (S. 20). \\[2pt]
20 & 1173 & \sichB{} & $f \in AC$ & $f \in \mathrm{AC}$ & Notation & Z. 1187 $\mathrm{AC}$; ebenso Z. 1362. \\[2pt]
20 & 1177 & \sichC{} & {quotient space \mfrac{F}{I}, where} & {quotient space $\mfrac{F}{I}$, where} & Mathemodus & \code{\textbackslash{}mfrac} ist für den Mathemodus gedacht (vgl. Z. 1088). \\[2pt]
21 & 1265 & \sichB{} & $\lim\sup_{t \to 0}$ & $\limsup_{t \to 0}$ & LaTeX & Index hängt sonst nur am sup; ebenso Z. 1512. \\[2pt]
22 & 1332 & \sichA{} & $\vert f_f'(0) - f_f'(s)\vert$ & $\vert \xi_f'(0) - \xi_f'(s)\vert$ & Tippfehler & Gemeint sind die Funktionen $\ensuremath{\xi}_f$ (Z. 1295, 1328); ebenso Z. 1337 $f'_f(0)$. \\[2pt]
22 & 1347 & \sichC{} & $0 \leq s \leq t_{0}$ & $0 < s \leq t_{0}$ & Tippfehler & Bei $s = 0$ ist 1/s nicht definiert; Z. 1332: $0 < s < t\ensuremath{_{0}}$. \\[2pt]
23 & 1398 & \sichA{} & {\emph{n-th Sobolev space}} & {\emph{$n$-th Sobolev space}} & Mathemodus & Zweimal in der Zeile (auch „n-th Sobolev semigroup“). \\[2pt]
26 & 1500 & \sichC{} & $E_{\mathcal{F}}^{T}$ & $E_{\mathcal{F}}^{\TT}$ & Notation & Daneben $m^{\TT}(E)$; ebenso Z. 1506, 1507, 1526, 1542 – bitte prüfen, ob gewollt. \\[2pt]
26 & 1522 & \sichB{} & $T(t)_{F}$ & $T(t)_{\mathcal{F}}$ & Notation & Gemeint ist der Filter $\ensuremath{\mathcal{F}}$, nicht der Raum F; ebenso Z. 1523. \\[2pt]
27 & 1583 & \sichC{} & {let $X$,$Y$ be} & {let $X$, $Y$ be} & Abstand & Leerzeichen nach dem Komma fehlt. \\[2pt]
27 & 1605 & \sichA{} & {on $E \times F$} & {on $E \otimes F$} & Tippfehler & Gemeint ist die Norm auf dem Tensorprodukt (Z. 1601, 1609). \\[2pt]
27 & 1628 & \sichC{} & $L^{1}(\mu \otimes \nu)$ & $L^{1}(\mu \times \nu)$ & Notation & Z. 1574, 1580: $\mu \times \nu$; ebenso Z. 1630. \\[2pt]
28 & 1683 & \sichB{} & $(T(h) \otimes S(h))(f \otimes g)$ & $(S(h) \otimes T(h))(f \otimes g)$ & Tippfehler & Im Satz heißt das Produkt $S(t)\ensuremath{\otimes}T(t)$; im Beweis vertauscht, auch Z. 1686. \\
\end{longtable}
\subsection*{Teil I, Kap. 2 – Characterization of Semigroups on Banach Spaces (\code{part-a/chap-a2.tex})}
\begin{longtable}{@{}rrc>{\RaggedRight\arraybackslash}p{5.8cm}>{\RaggedRight\arraybackslash}p{5.8cm}>{\RaggedRight\arraybackslash}p{1.9cm}>{\RaggedRight\arraybackslash}p{7.3cm}@{}}
\toprule
\textbf{S.} & \textbf{Z.} & \textbf{} & \textbf{Ist} & \textbf{Soll} & \textbf{Art} & \textbf{Kommentar} \\
\midrule
\endfirsthead
\toprule
\textbf{S.} & \textbf{Z.} & \textbf{} & \textbf{Ist} & \textbf{Soll} & \textbf{Art} & \textbf{Kommentar} \\
\midrule
\endhead
\bottomrule
\endlastfoot
32 & 92 & \sichB{} & $u'(t) = Au(t)$ & $\dot{u}(t) = Au(t)$ & Notation & Sonst durchgehend $\dot{u}$ (z. B. Z. 86). \\[2pt]
33 & 148 & \sichA{} & {$\tilde{v}(s) = v(s)$ \& \textbackslash{}text\{$0 \leq s \leq t$.\}} & {v(s) \& \textbackslash{}text\{for $0 \leq s \leq t$.\}} & Tippfehler & „\ensuremath{\tilde{v}}(s) $=$“ im zweiten cases-Zweig doppelt; „for“ fehlt. \\[2pt]
36 & 388 & \sichC{} & {Then $ \overline{A} $ satisfies} & {Then $ A $ satisfies} & Tippfehler & Eingeführt wird die Einschränkung A, nicht ihr Abschluss; bitte prüfen. \\[2pt]
36 & 417 & \sichC{} & $u(t) = (f + tBg,g)$ & $u(t) = (f - tBg,g)$ & Tippfehler & Mit $A =$ ($0 \ensuremath{-}B$; 0 0) wäre \ensuremath{\dot{u}} $=$ ($\ensuremath{-}Bg$, 0); nur Hinweis. \\[2pt]
38 & 528 & \sichA{} & {is dense in E and} & {is dense in $E$ and} & Mathemodus & E ohne Mathemodus. \\[2pt]
39 & 552 & \sichC{} & {~\lbl{cor:a2-1.8.2} implies~\lbl{cor:a2-1.8.1}} & {~\lbl{thm:a2-1.7.2} implies~\lbl{thm:a2-1.7.1}} & Verweis & Gedruckt gleich, aber der Hyperlink führt ins Korollar statt ins Theorem. \\[2pt]
40 & 657 & \sichA{} & $(\lambda \in S(\alpha_{1} + \pi/2).$ & $(\lambda \in S(\alpha_{1} + \pi/2)).$ & Klammern & Äußere Klammer nicht geschlossen. \\[2pt]
40 & 686 & \sichB{} & $\lambda \in \rho(A)$ & $\lambda \in \rho(A^{2})$ & Tippfehler & Gemeint ist die Resolventenmenge von $A\ensuremath{^{2}}$ (vgl. Z. 957). \\[2pt]
41 & 744 & \sichB{} & {$(U_{i})_{t \geq 0}$ the group} & {$(U_{i}(t))_{t \in \R}$ the group} & Tippfehler & (t) fehlt; Gruppe, daher $t \ensuremath{\in} \ensuremath{\mathbb{R}}$ (vgl. Z. 747). \\[2pt]
42 & 853 & \sichB{} & $\alpha \in \left]0,\pi/2\right]$ & $\alpha \in (0,\pi/2]$ & Notation & Sonst $(0,\pi/2]$ (Z. 604, 647, 798); ebenso Z. 1417 $\left]0,1\right]$. \\[2pt]
43 & 940 & \sichB{} & $\Im\{z(t)\}^2/\Re\{z(t)\}^2$ & $\Im(z(t)^2)/\Re(z(t)^2)$ & Klammern & Geschweifte Klammern werden gedruckt und lesen sich als (Im $z)\ensuremath{^{2}}$; gemeint $\Im (z\ensuremath{^{2}})$ (Z. 934). \\[2pt]
43 & 949 & \sichA{} & $(r\eu^{\im\theta}+\omega)^2$ & $(r\eu^{\im\theta}+w)^2$ & Tippfehler & Im Beweis heißt es sonst w (Z. 902, 954). \\[2pt]
45 & 1020 & \sichC{} & $c\eu^{-b\Re(\lambda)}$ & $c\eu^{-b\cdot\Re(\lambda)}$ & Notation & Theorem 1.17 und 1.18 (Z. 1036, 1043) einmal mit, einmal ohne $\cdot$. \\[2pt]
46 & 1108 & \sichC{} & {[Proof of Theorem~\lbl{thm:a2-1.20}  ]} & {[Proof of Theorem~\lbl{thm:a2-1.20}]} & LaTeX & Schon in der Arbeitsliste (S. 46); ebenso Z. 2502. \\[2pt]
46 & 1126 & \sichA{} & $T(t) \dt (n \in \N).$ & $T(t) \dt \quad (n \in \N).$ & Abstand & Gedruckt „$dt(n \ensuremath{\in} \ensuremath{\mathbb{N}}$)“. \\[2pt]
46 & 1152 & \sichC{} & \code{\textbackslash{}right\allowbreak{}\textbackslash{}lVert} & \code{\textbackslash{}right\allowbreak{}\textbackslash{}rVert} & LaTeX & Falsches schließendes Normzeichen (optisch gleich). \\[2pt]
46 & 1173 & \sichC{} & $T(t)\,\dt$ & $T(t)\dt$ & Abstand & Zusätzliches \code{\textbackslash{},} nur in Z. 1173–1180, 1544, 1551, 1638, 1644–1647. \\[2pt]
50 & 1370 & \sichB{} & $[\text{ess}]\text{-}\text{image}(m)$ & wie Z. 1455 ($\Image{m}$) & Notation & Dasselbe Objekt einmal „image“, einmal „im“. \\[2pt]
50 & 1411 & \sichB{} & $H \coloneq \{m(x) \colon \Re(\lambda) \geq b \dots$ & $H \coloneq \{\lambda \in \C \colon \Re(\lambda) \geq b \dots$ & Tippfehler & Bedingung über $\ensuremath{\lambda}$, im Mengenterm steht m(x). \\[2pt]
50 & 1420 & \sichA{} & \code{[\allowbreak{}\textbackslash{}text\{ess\}]\allowbreak{}\textbackslash{}text\{-\}\allowbreak{}\textbackslash{}esssup\allowbreak{}\textbackslash{}left\allowbreak{}\textbackslash{}\{} & \code{[\allowbreak{}\textbackslash{}text\{ess\}]\allowbreak{}\textbackslash{}text\{-\}\allowbreak{}\textbackslash{}sup\allowbreak{}\textbackslash{}left\allowbreak{}\textbackslash{}\{} & Notation & Gedruckt „[ess]-ess sup“; vgl. Z. 1427. \\[2pt]
50 & 1473 & \sichA{} & $S(\theta + \pi/2) \subset~\rho(A)$ & $S(\theta + \pi/2) \subseteq \rho(A)$ & LaTeX & Tilde im Mathemodus erzeugt Leerraum; sonst $\subseteq$. \\[2pt]
50 & 1509 & \sichC{} & $\iu \cdot \exp(n^{2})$ & $\iu\,\eu^{n^{2}}$ & Notation & Die drei Beispiele (Z. 1509, 1513, 1517) uneinheitlich. \\[2pt]
53 & 1628 & \sichB{} & $(T(t+s-h) - T(t-s))$ & $(T(t+h-s) - T(t-s))$ & Tippfehler & Z. 1647 hat $t+h\ensuremath{-}s$; ebenso Z. 1634, 1638. \\[2pt]
54 & 1645 & \sichA{} & $\int_{0}^{t} (T(t-s)BS(s)\, \ds$ & $\int_{0}^{t} T(t-s)BS(s)\, \ds$ & Klammern & Öffnende Klammer ohne schließende. \\[2pt]
54 & 1646 & \sichC{} & $\leq \|T(t+h) - T(t)\|  + \dots$ & $\leq \|T(t+h) - T(t)\| + \dots$ & LaTeX & $\ensuremath{\leq}$ links vom Ausrichtungspunkt, Ausrichtung verschoben. \\[2pt]
54 & 1675 & \sichA{} & $= (\Id - SR(\lambda,A)$ & $= (\Id - SR(\lambda,A))$ & Klammern & Schließende Klammer fehlt. \\[2pt]
54 & 1679 & \sichA{} & $(\Id -BR(\lambda,A))^{-1})$ & $(\Id -BR(\lambda,A))^{-1}$ & Klammern & Überzählige schließende Klammer. \\[2pt]
57 & 1819 & \sichC{} & {\emph{p-contractive}}, {p-dissipative} & {$p$-contractive}, {$p$-dissipative} & Mathemodus & p steht aufrecht, {$N$-}, {$p$-} (Z. 2248) dagegen kursiv; ca. 30 Stellen in Z. 1819–2272. \\[2pt]
58 & 1914 & \sichA{} & $\phi~\in~\ddp(f)$ & $\phi \in \ddp(f)$ & LaTeX & Tilden im Mathemodus erzeugen Leerraum. \\[2pt]
58 & 1947 & \sichC{} & {for all $t > 0$, \code{\textbackslash\textbackslash}} & {for all $t > 0$,} & LaTeX & \code{\textbackslash{}\allowbreak{}\textbackslash{}} vor abgesetzter Formel. \\[2pt]
58 & 1951 & \sichA{} & $p(f)~  \leq~p(f-tAf)$ & $p(f)  \leq p(f-tAf)$ & LaTeX & Tilden im Mathemodus. \\[2pt]
59 & 2031 & \sichA{} & \code{T(t) \allowbreak{}= \allowbreak{}\allowbreak{}\textbackslash{}lim\_\{t\allowbreak{}\textbackslash{}to \allowbreak{}0\} \allowbreak{}\allowbreak{}\textbackslash{}left( \allowbreak{}\allowbreak{}\textbackslash{}frac\{n\}\{t\} \allowbreak{}\dots{}} & \code{T(t) \allowbreak{}= \allowbreak{}\allowbreak{}\textbackslash{}lim\_\{n\allowbreak{}\textbackslash{}to\allowbreak{}\textbackslash{}infty\} \allowbreak{}\allowbreak{}\textbackslash{}left( \allowbreak{}\allowbreak{}\textbackslash{}frac\{n\}\{t\} \allowbreak{}\dots{}} & Tippfehler & Grenzübergang über n (vgl. Z. 584). \\[2pt]
60 & 2116 & \sichC{} & {(this is equivalent \dots  {schaefer:1966}).} & {\dots  {schaefer:1966})).} & Klammern & Schon in der Arbeitsliste (S. 61). \\[2pt]
62 & 2235 & \sichB{} & {from~\lbl{thm:a2-2.6} it follows} & {from Theorem~\lbl{thm:a2-2.6} it follows} & Verweis & Schon in der Arbeitsliste (S. 62). \\[2pt]
63 & 2274 & \sichB{} & $p((1-tA)f) \geq p(f)$ & $p((\Id-tA)f) \geq p(f)$ & Notation & Identität sonst $\Id$. \\[2pt]
63 & 2281 & \sichA{} & {if it is N-dissipative} & {if it is $N$-dissipative} & Mathemodus & Z. 1861 schreibt {$N$-dissipative}. \\[2pt]
64 & 2311 & \sichB{} & $\lim\sup_{t\to 0}$ & $\limsup_{t\to 0}$ & Notation & Z. 2014 $\limsup$; ebenso Z. 2322, 2323, 2341, 2431. \\[2pt]
64 & 2341 & \sichC{} & $< 1\ ,$ & $< 1 ,$ & Satzzeichen & Leerzeichen vor dem Komma. \\[2pt]
66 & 2513 & \sichB{} & {and $H^{\infty}$ does not only contain} & {and $H^{\infty}(O)$ does not only contain} & Tippfehler & Im selben Satz $H^{\infty}(O)$. \\
\end{longtable}
\subsection*{Teil I, Kap. 3 – Spectral Theory (\code{part-a/chap-a3.tex})}
\begin{longtable}{@{}rrc>{\RaggedRight\arraybackslash}p{5.8cm}>{\RaggedRight\arraybackslash}p{5.8cm}>{\RaggedRight\arraybackslash}p{1.9cm}>{\RaggedRight\arraybackslash}p{7.3cm}@{}}
\toprule
\textbf{S.} & \textbf{Z.} & \textbf{} & \textbf{Ist} & \textbf{Soll} & \textbf{Art} & \textbf{Kommentar} \\
\midrule
\endfirsthead
\toprule
\textbf{S.} & \textbf{Z.} & \textbf{} & \textbf{Ist} & \textbf{Soll} & \textbf{Art} & \textbf{Kommentar} \\
\midrule
\endhead
\bottomrule
\endlastfoot
73 & 182 & \sichA{} & \code{\textbackslash{}begin\{example\}\allowbreak{}\textbackslash{}label\{ex:a3-1.4\}[Hilbert \allowbreak{}space, \allowbreak{}\allowbreak{}\textbackslash{}textcite\{zabczyk:1975\}]} & \code{\textbackslash{}begin\{example\}[Hilbert \allowbreak{}space, \allowbreak{}\allowbreak{}\textbackslash{}textcite\{zabczyk:1975\}]\allowbreak{}\textbackslash{}label\{ex:a3-1.4\}} & LaTeX & Optionales Argument hinter \code{\textbackslash{}label}: „[Hilbert space, \dots{}]“ erscheint als Fließtext statt im Kopf des Beispiels (vgl. Z. 127). \\[2pt]
73 & 207 & \sichC{} & {$n \in \N$ ,} & {$n \in \N$,} & Satzzeichen & Leerzeichen vor dem Komma; dasselbe Muster in Z. 210, 712, 1050, 1053, 1056, 1162, 1163, 1166, 1215, 1383, 1386, 1432, 1490, 1493, 1713, 1733, 1736, 2206. \\[2pt]
73 & 234 & \sichB{} & {continuous on each component $E_{n}$} & $H_{n}$ & Tippfehler & Komponenten heißen $H_n$ (Z. 184, 220). \\[2pt]
76 & 355 & \sichA{} & $(\text{dist}(x,\sigma(A)))^{-1}$ & $(\mathrm{dist}(\lambda_{n},\sigma(A)))^{-1}$ & Tippfehler & x statt $\ensuremath{\lambda}_n$; sonst $\mathrm{dist}$ (Z. 537). \\[2pt]
76 & 358 & \sichB{} & {Proposition~\lbl{prop:a3-2.5}\,(ii)} & $\dots \,(i)$ & Verweis & Schon in der Arbeitsliste (S. 76). \\[2pt]
76 & 392 & \sichB{} & $q \colon X \mapsto \C$ & $q \colon X \to \C$ & Notation & Sonst $\to$ (Z. 1933, 1948). \\[2pt]
76 & 402 & \sichC{} & $\eu^{tq} \cdot f \, , \, f \in E .$ & $\eu^{tq} \cdot f, \quad f \in E .$ & Abstand & Sonst $, \quad$; ebenso Z. 2141, 2293. \\[2pt]
76 & 431 & \sichC{} & {spectral value of\/ $A$} \dots{} & $M_{q}$ & Notation & Im Beispiel heißt der Generator sonst $M_q$ (Z. 431–434). \\[2pt]
77 & 467 & \sichC{} & $\epsilon_{\lambda}$ (vs. $e_{\lambda}$ Z. 160) & einheitlich $\epsilon_{\lambda}$ & Notation & Sonst $\ensuremath{\varepsilon}_\ensuremath{\lambda}$ (Z. 520, 708, 1515). \\[2pt]
78 & 603 & \sichB{} & $\diff{z} =    = \dots$ & $\diff{z}    = \dots$ & Satzzeichen & Doppeltes Gleichheitszeichen am Zeilenende und -anfang. \\[2pt]
82 & 757 & \sichB{} & $T_{1} = (T(t)P)_{t \geq 0}$ & $\TT_{1} = \dots$ & Notation & Halbgruppen sonst $\ensuremath{\mathcal{T}}_i$ (Z. 669, 685, 727); ebenso Z. 774 $(T\ensuremath{_{3}})$, 801 $(\ensuremath{\omega}\ensuremath{_{0}}(T\ensuremath{_{1}}))$. \\[2pt]
82 & 808 & \sichA{} & {Theorem~~\lbl{thm:a3-3.3}} & {Theorem~\lbl{thm:a3-3.3}} & LaTeX & Doppelte Tilde, sichtbar doppelter Abstand. \\[2pt]
83 & 817 & \sichB{} & {(see~\lbl{prop:a3-6.2})} & {(see Proposition~\lbl{prop:a3-6.2})} & Verweis & Schon in der Arbeitsliste (S. 83). \\[2pt]
86 & 1070 & \sichB{} & $D(A)\backslash N$ & $D(A) \setminus N$ & Notation & Sonst $\setminus$ (Z. 537, 1446). \\[2pt]
86 & 1071 & \sichC{} & $E_{/}$ & $E/N$ & Notation & Z. 1018: E/N; ebenso Z. 1077, 1084. (Fehlender Schlusspunkt schon in der Arbeitsliste, S. 87.) \\[2pt]
87 & 1134 & \sichC{} & {we have \quad $\dots $, $\quad R\dots  = \sum U_{n\vert}\dots $} & abgesetzte Formel, einheitlich $\sum_{n}$ & LaTeX & $\quad$ im Fließtext, Tilde am Ende, Summe teils ohne Index (Z. 1134–1135). \\[2pt]
88 & 1155 & \sichA{} & $\lim_{t \to \infty} \|T(t)'\phi-\phi\| = 0$ & $\lim_{t \to 0} \dots$ & Tippfehler & Definition des Halbgruppen-Duals über $t \ensuremath{\to} 0$ (vgl. chap-a1 Z. 1245). \\[2pt]
88 & 1173 & \sichB{} & {\lbl{prop:a3-2.2}(ii)} & {Proposition~\lbl{prop:a3-2.2}\,(ii)} & Verweis & Schon in der Arbeitsliste (S. 88). \\[2pt]
89 & 1233 & \sichA{} & $({\hat{f}^{m}})_{n \in \N}$ & $(\hat{f}^{m})_{m \in \N}$ & Tippfehler & Folge läuft über m. \\[2pt]
91 & 1351 & \sichA{} & {the n-th Fourier coefficient} & {the $n$-th Fourier coefficient} & Mathemodus & n nicht im Mathemodus. \\[2pt]
93 & 1511 & \sichC{} & {generator\quad $Af(z) = \dots $\quad is \dots } & ohne $\quad$ & Abstand & $\quad$ im Fließtext (Z. 1511–1514). \\[2pt]
93 & 1525 & \sichC{} & $Q_{m} \coloneq \sum_{-m}^{+m} P_{n}$ in $L^{p}(\Gamma,m)$ & anderer Index, z. B. $Q_{N}$ & Notation & m zugleich Summationsgrenze und Maß. \\[2pt]
94 & 1574 & \sichC{} & $\sigma(T(t)) = \Gamma$ & $\sigma(R_{\tau}(t)) = \Gamma$ & Notation & In (ii) heißt die Gruppe $R_\ensuremath{\tau}(t)$ (Z. 1572). \\[2pt]
96 & 1702 & \sichB{} & {(see Proposition~\lbl{prop:a3-4.4})} & {(see Section~\lbl{subsec:a3-4.4})} & Verweis & Schon in der Arbeitsliste (S. 96). \\[2pt]
97 & 1736 & \sichB{} & $\overline{\text{span}_{n \in \Z} \{P_{n}E\}}$ & $\overline{\mathrm{span}_{n \in \Z} \dots }$ & Notation & $\text{span}$ wird im kursiven Korollar kursiv; Z. 1716 $\mathrm{span}$. \\[2pt]
97 & 1772 & \sichC{} & \code{\textbackslash{}label\{eq:a3-6.6)\}} & \code{\textbackslash{}label\{eq:a3-6.6\}} & LaTeX & Überzählige Klammer im Label (derzeit nicht referenziert). \\[2pt]
97 & 1815 & \sichB{} & {(see~\lbl{cor:a3-1.2})} & {(see Corollary~\lbl{cor:a3-1.2})} & Verweis & Bezeichnung fehlt, gedruckt nur die Nummer. \\[2pt]
98 & 1898 & \sichC{} & {Theorem~\lbl{thm:a2-1.20} )} & {Theorem~\lbl{thm:a2-1.20})} & Satzzeichen & Schon in der Arbeitsliste (S. 99). \\[2pt]
102 & 2099 & \sichB{} & $\widehat{\left(\eu^{\im 2k}\check{\phi_{0}}\right)}$ & $\widehat{\left(\eu^{\im 2k\cdot}\check{\phi_{0}}\right)}$ & Tippfehler & Variable fehlt im Exponenten (vgl. Z. 2161 $\eu^{\im\pi\cdot}$). \\[2pt]
103 & 2150 & \sichC{} & $\eu^{-\im n0}$ & $\eu^{-\im n\cdot 0}$ & Notation & „n0“ liest sich wie ein Symbol. \\[2pt]
105 & 2194 & \sichA{} & $C_{1}\text{-}\sum_{k\in\Z} \xi_{n}$ & $C_{1}\text{-}\sum_{k\in\Z} \xi_{k}$ & Tippfehler & Summationsindex k, Summand mit n. \\[2pt]
106 & 2269 & \sichB{} & \code{< \allowbreak{}\allowbreak{}\textbackslash{}infty\$\textasciitilde{}, \allowbreak{}then\allowbreak{}\textbackslash{}quad \allowbreak{}\$\dots{}\$\allowbreak{}\textbackslash{}quad \allowbreak{}is \allowbreak{}summable} & \code{< \allowbreak{}\allowbreak{}\textbackslash{}infty\$, \allowbreak{}then \allowbreak{}\$\dots{}\$ \allowbreak{}is \allowbreak{}summable} & Satzzeichen & Geschütztes Leerzeichen vor Komma, $\quad$ im Fließtext. \\[2pt]
106 & 2270 & \sichC{} & Leerzeile in Item (ii) & Leerzeile entfernen & LaTeX & Absatzumbruch mitten im Satz. \\[2pt]
106 & 2286 & \sichB{} & {(use Eq.~(\lbl{eq:a3-7.5}))} & \code{eq:a3-7.7} & Verweis & Benötigt wird die Definition von $Q_k$ (7.7), nicht (7.5). \\[2pt]
106 & 2311 & \sichA{} & \code{\textbackslash{}eu\textasciicircum{}\{-\allowbreak{}\textbackslash{}im \allowbreak{}ks\}g\_\{k\}\allowbreak{}\textbackslash{}right\allowbreak{}\textbackslash{}|\textasciicircum{}\{2\} \allowbreak{}\allowbreak{}\textbackslash{}dx} & $\dots  \ds$ & Tippfehler & Integrationsvariable ist s. \\[2pt]
106 & 2313 & \sichC{} & $c(\sum_{k\in N} \|g_{k}\|^{2})^{1/2}$ & $c\left(\sum \dots  \right)^{1/2}$ & Klammern & Klammern zu klein; daneben mit $\left(\dots \right)$. \\[2pt]
107 & 2335 & \sichC{} & $\not\in$ $(2\ensuremath{\times})$ & $\notin$ & Notation & Sonst $\notin$. \\
\end{longtable}
\subsection*{Teil I, Kap. 4 – Asymptotics of Semigroups on Banach Spaces (\code{part-a/chap-a4.tex})}
\begin{longtable}{@{}rrc>{\RaggedRight\arraybackslash}p{5.8cm}>{\RaggedRight\arraybackslash}p{5.8cm}>{\RaggedRight\arraybackslash}p{1.9cm}>{\RaggedRight\arraybackslash}p{7.3cm}@{}}
\toprule
\textbf{S.} & \textbf{Z.} & \textbf{} & \textbf{Ist} & \textbf{Soll} & \textbf{Art} & \textbf{Kommentar} \\
\midrule
\endfirsthead
\toprule
\textbf{S.} & \textbf{Z.} & \textbf{} & \textbf{Ist} & \textbf{Soll} & \textbf{Art} & \textbf{Kommentar} \\
\midrule
\endhead
\bottomrule
\endlastfoot
111 & 54 & \sichA{} & {$ f \in \C^{n} $ ,} & {$ f \in \C^{n} $,} & Satzzeichen & Leerzeichen vor dem Komma; ebenso Z. 693. \\[2pt]
111 & 71 & \sichB{} & {Chapter~\lbl{chap:c3}} & {Chapter~\lbl{chap:a3}} & Verweis & Gemeint ist wohl Kap. 3 (Spektraltheorie) dieses Teils; Verweise in andere Teile haben sonst „Part-\dots{}“. \\[2pt]
112 & 79 & \sichB{} & {($1^{st}$ part)} & $(1\textsuperscript{st} part)$ & Mathemodus & „st“ kursiv; ebenso Z. 111 $2^{nd}$. \\[2pt]
113 & 176 & \sichA{} & $\Im(\Delta) \not= C_0(\R^{n})$ & $\Image{\Delta} \not= C_0(\R^{n})$ & Notation & $\Im$ ist der Imaginärteil; Bild sonst $\Image$ (Z. 888). \\[2pt]
113 & 205 & \sichA{} & $\int_{0}^{\infty} \eu^{\lambda t} T(t)f \dt$ & $\int_{0}^{\infty} \eu^{-\lambda t} T(t)f \dt$ & Tippfehler & Minus im Exponenten fehlt (vgl. Z. 528). \\[2pt]
113 & 219 & \sichC{} & $\|T(t)f\| \leq M\eu^{\lambda t}$ & $\|T(t)f\| \leq M\eu^{\Re(\lambda) t}$ & Notation & $\ensuremath{\lambda}$ ist komplex (Z. 205). \\[2pt]
115 & 255 & \sichB{} & $\int_{0}^{t} \eu^{\mu s} f(s) \ds$ & $\int_{0}^{t} \eu^{-\mu s} f(s) \ds$ & Tippfehler & Ebenso Z. 256 $\eu^{\mu r}$; passt dann zu Z. 262. \\[2pt]
115 & 292 & \sichB{} & \code{\textbackslash{}begin\{enumerate\}[(i)]} & \code{\textbackslash{}begin\{enumerate\}[\allowbreak{}\textbackslash{}upshape \allowbreak{}(i)]} & Notation & Labels werden sonst kursiv; ebenso Z. 300, 910. \\[2pt]
115 & 295 & \sichC{} & $\inf\{\Re(\lambda)\colon\int_{0}^{\infty} \dots \}$ & $\inf\left\{\dots \right\}$ & Klammern & Kleine Mengenklammern um Integral; ebenso Z. 303, 357 (Z. 339, 465, 470 skalieren). \\[2pt]
116 & 339 & \sichA{} & $\|^{-\lambda t} T(t)f\|$ & $\|\eu^{-\lambda t} T(t)f\|$ & Tippfehler & $\eu$ fehlt. \\[2pt]
116 & 361 & \sichB{} & {from Eq.~(\lbl{eq:a4-1})} & {from Eq.~(\lbl{eq:a4-1.1})} & Verweis & eq:a4-1 ist (*); gemeint (1.1); ebenso Z. 457. \\[2pt]
116 & 362 & \sichC{} & {that \[g \coloneqq\lim_{r \to \infty} T(r)f\] exists.} & {that $g \coloneq \lim_{r \to \infty} T(r)f$ exists.} & LaTeX & Abgesetzte Formel mitten im Satz; sonst $\coloneq$. \\[2pt]
118 & 474 & \sichC{} & $\eu^{\lambda t} \left\{ f - \int \dots  \right\}$ & $\eu^{\lambda t} \left( f - \int \dots  \right)$ & Notation & Z. 209 mit runden Klammern. \\[2pt]
119 & 552 & \sichB{} & $x^n \in \C^n\colon\sum_{j=1}^{\infty} \dots$ & $x^n \in \C^n,\ \sum_{j=1}^{\infty} \dots$ & Notation & Zwei $\colon$ in einer Mengenklammer; $...$ statt $\dots$ (auch Z. 569, 588). \\[2pt]
119 & 554 & \sichB{} & $\eu^{2\pi i n t}$ & $\eu^{2\pi \im n t}$ & Notation & Kursives i; sonst $\im$; ebenso Z. 567, 1167. \\[2pt]
119 & 574 & \sichA{} & $\|\exp( t A_n) \cdot e_n\|^2 =    = \dots$ & $\dots  \|^2    = \dots$ & Tippfehler & Doppeltes Gleichheitszeichen. \\[2pt]
119 & 588 & \sichC{} & $x_q :=  (q e_1, \dots )$ & $x_q \coloneq (q e_1, \dots )$ & Notation & Sonst $\coloneq$. \\[2pt]
119 & 593 & \sichB{} & $\| \exp(t A) e_n \|^2$ & $\| \exp(t A_n) e_n \|^2$ & Tippfehler & $e_n \ensuremath{\in} \ensuremath{\mathbb{C}}\ensuremath{^{n}}$, gemeint ist $A_n$ (Z. 574). \\[2pt]
121 & 633 & \sichB{} & {\lbl{cor:a3-6.7} ).} & {\lbl{cor:a3-6.7}).} & LaTeX & Teilweise schon in der Arbeitsliste (S. 121–125); ebenso Z. 638, 792, 837, 879. \\[2pt]
121 & 634 & \sichC{} & $\dot{u}(t) = A u(t), u(0) = x$ & $\dot{u}(t) = A u(t),\ u(0) = f$ & Notation & Anfangswert sonst f (Z. 41). \\[2pt]
121 & 682 & \sichB{} & $\sup \{ \omega(f),  f \in D(A^n) \}$ & $\sup \{ \omega(f) \colon f \in D(A^n) \}$ & Notation & Sonst $\colon$. \\[2pt]
121 & 690 & \sichC{} & $\max(\omega_{0}(A), q)$ & $\max\{\omega_{0}(A), q\}$ & Notation & Z. 692, 613, 797: $\max\{\dots \}$. \\[2pt]
122 & 695 & \sichB{} & $\frac{1}{2\pi i} \dots  f \, d\lambda$ & $\frac{1}{2\pi\im} \dots  f \diff{\lambda}$ & Notation & Kursives i und d; ebenso Z. 707, 709, 718, 1271. \\[2pt]
122 & 711 & \sichB{} & Leerzeile nach \code{\textbackslash{}end\{equation\}} & Leerzeile löschen & LaTeX & Fortsetzungssatz bekommt sonst Absatzeinzug. \\[2pt]
123 & 755 & \sichB{} & $(a)  \Rightarrow  (b)  \Rightarrow  (c)$ & {(a)~$\Rightarrow$ (b)~$\Rightarrow$ (c)} & Notation & Labels im Mathemodus kursiv; ebenso Z. 766, 923, 928, 930, 935, 937. \\[2pt]
124 & 840 & \sichB{} & $T(t)^* \phi = \exp(\lambda t) \phi$ & $T(t)' \phi = \eu^{\lambda t} \phi$ & Notation & Adjungierte sonst mit $\ensuremath{{}'}$ (Z. 839); Z. 837 $\eu^{\lambda t}$. \\[2pt]
127 & 1015 & \sichB{} & $F(t) \quad , \quad u(0) = f,$ & $F(t), \quad u(0) = f,$ & Satzzeichen & Großer Abstand vor dem Komma. \\[2pt]
127 & 1050 & \sichB{} & $t\to T(t)g$ & $t\mapsto T(t)g$ & Notation & Abbildungsvorschrift (vgl. Z. 1044). \\[2pt]
128 & 1095 & \sichB{} & {Theorem~\lbl{prop:a1-1.6}} & {Proposition~\lbl{prop:a1-1.6}} & Verweis & Label gehört zu einer Proposition. \\[2pt]
130 & 1202 & \sichB{} & $\{f \in E: f' \in C^{1} \text{ and } f' \in E\}$ & $\{f \in E\colon f \in C^{1} \text{ and } f' \in E\}$ & Tippfehler & Gemeint $f \ensuremath{\in} C\ensuremath{^{1}}$; $:$ statt $\colon$. \\[2pt]
130 & 1208 & \sichC{} & \code{\textbackslash{}tag\{*\}} & anderes Tag, z. B. {\textbackslash{}tag\{$\ast\ast\ast$\}} & Verweis & (*) schon in Z. 24 vergeben; Bezug in Z. 1212/1234 mehrdeutig. Kursives d in $\frac{d}{dx}$ (Z. 1202, 1208). \\[2pt]
130 & 1254 & \sichC{} & {Example~\lbl{ex:a4-1.2}(ii)} & {Example~\lbl{ex:a4-1.2}\,(ii)} & Abstand & Sonst mit \code{\textbackslash{},}. \\[2pt]
131 & 1277 & \sichC{} & {statements Eq.~(\lbl{eq:a4-1.12}), Eq.~(\lbl{eq:a4-1.13})} & prüfen & Verweis & (1.12), (1.13) sind Inversionsformeln; vermutlich Nummerierung der 1. Auflage. \\[2pt]
131 & 1279 & \sichC{} & {[\lbl{datko:1970}]} & prüfen & Verweis & Z. 802 nennt \code{datko:1972}. \\
\end{longtable}
\end{document}
